%%%%%%%%%%  scan idx 72  |  folio 59  %%%%%%%%%%
$X\times k$ ($n$ est la
dimension de $X$), et deux points $a,a'\in k$ tels que l'on ait~:
\[
 Z=f_a^{-1}(Y),\qquad Z'=f_{a'}^{-1}(Y),
\]
où l'on pose $f_u(x)=(x,u)$. En vertu de la formule (2.23), on a alors~:
\[
 \gamma_X(Z)\equiv f_a^!(\gamma_{X\times k}(Y))\qquad\mathrm{mod}\ \underline K(X)^{p+1}
\]
et une formule analogue pour $\gamma_X(Z')$. Il suffit alors de prouver que $f_a^!=f_{a'}^!$, ce qui résulte aussitôt de la proposition 2.9.
\begin{flushright}cqfd.\end{flushright}

\textbf{Corollaire 1~:} \emph{La filtration de $\underline K(X)$ est compatible avec sa structure d'anneau, i.e.\ on a $\underline K(X)^i\underline K(X)^j\subset\underline K(X)^{i+j}$ quels que soient les entiers $i,j\geq0$. On a de plus un homomorphisme surjectif naturel d'anneaux gradués~:}
\begin{equation}\tag{2.34}
 \varphi:A(X)\longrightarrow G(\underline K(X))
\end{equation}
(où le deuxième membre désigne l'anneau gradué associé à $\underline K(X)$ muni de la filtration $(\underline K(X)^i)_{i\geq0}$) \emph{défini par la condition}~:
\[
 \varphi(Z^{n-p})=\gamma(Z^{n-p})\qquad\mathrm{mod}\ \underline K(X)^{p+1}.
\]
Ce corollaire résulte immédiatement du théorème 2.12, des propositions 2.6 et 2.7 et du théorème de Chow 2.1, par récurrence sur l'entier $i+j$.

\textbf{Corollaire 2~:} \emph{Soient $X$ et $Y$ deux variétés quasi-projectives non singulières et $f$ un morphisme de $X$ dans $Y$. Alors l'homomorphisme $f^!$ de $\underline K(Y)$ dans $\underline K(X)$ est compatible avec les filtrations, et l'on a un diagramme commutatif~:}
\[
\begin{tikzcd}[column sep=large,row sep=large]
G(\underline K(Y)) \arrow[r,"G(f^!)"] & G(\underline K(X))\\
A(Y) \arrow[r,"f^*"] \arrow[u,"\varphi_Y"'] & A(X) \arrow[u,"\varphi_X"'] .
\end{tikzcd}
\]
